\documentclass[12pt,a4paper,oneside]{report}

\usepackage[utf8]{inputenc}
\usepackage[T1]{fontenc}
\usepackage[portuguese]{babel}
\usepackage{lmodern}
\usepackage{geometry}
\usepackage{setspace}
\usepackage{graphicx}
\usepackage{float}
\usepackage{booktabs}
\usepackage{longtable}
\usepackage{array}
\usepackage{amsmath}
\usepackage{hyperref}
\usepackage{xcolor}
\usepackage{listings}
\usepackage{caption}
\usepackage{enumitem}
\usepackage{fancyhdr}
\usepackage{titlesec}
\usepackage{helvet}
\usepackage[nohyperlinks]{acronym}

\geometry{left=3cm,right=2.5cm,top=3cm,bottom=3cm}
\onehalfspacing
\setlength{\parindent}{1.25cm}
\setlength{\parskip}{0.2em}
\setlength{\headheight}{16pt}

\hypersetup{
    colorlinks=true,
    linkcolor=black,
    citecolor=black,
    urlcolor=blue,
    pdftitle={MindBridgeXR},
    pdfauthor={Hugo Serra}
}

\lstdefinelanguage{CSharp}{
    language=[Sharp]C
}
\lstdefinestyle{codestyle}{
    basicstyle=\ttfamily\small,
    breaklines=true,
    frame=single,
    numbers=left,
    numberstyle=\tiny,
    captionpos=b,
    showstringspaces=false,
    tabsize=4
}
\lstset{style=codestyle}

\newcommand{\blankpage}{
    \clearpage
    \thispagestyle{empty}
    \null
    \clearpage
}

\renewcommand{\chaptermark}[1]{\markboth{\MakeUppercase{#1}}{}}
\renewcommand{\sectionmark}[1]{}
\renewcommand{\bibname}{Bibliografia}

\pagestyle{fancy}
\fancyhf{}
\fancyhead[L]{\bfseries\thepage}
\fancyhead[R]{\bfseries\nouppercase{\leftmark}}
\renewcommand{\headrulewidth}{0.4pt}
\renewcommand{\footrulewidth}{0pt}

\fancypagestyle{plain}{
    \fancyhf{}
    \fancyhead[L]{\bfseries\thepage}
    \fancyhead[R]{\bfseries\nouppercase{\leftmark}}
    \renewcommand{\headrulewidth}{0.4pt}
    \renewcommand{\footrulewidth}{0pt}
}

\titleformat{\chapter}[display]
    {\normalfont}
    {\Large\bfseries\itshape\chaptertitlename\\[0.35cm]
     {\fontfamily{phv}\selectfont\fontsize{64}{72}\selectfont\bfseries\color{gray}\thechapter}}
    {1cm}
    {\Huge\bfseries}

\titlespacing*{\chapter}{0pt}{1.5cm}{2cm}

\begin{document}

\begin{titlepage}
    \centering

    {\Huge\bfseries Universidade da Beira Interior\par}
    {\Huge Departamento de Informática\par}

    \vspace{2cm}

    \includegraphics[width=0.32\textwidth]{ubi.png}

    \vspace{1.8cm}

    {\Large\bfseries\itshape MindBridgeXR\par}
    {\Large\bfseries\itshape Ambientes Imersivos para Estimulação Cognitiva em Alzheimer\par}

    \vspace{1.2cm}

    {\Large\bfseries Hugo Maia Serra\par}

    \vspace{1.8cm}

    {\Large\bfseries Licenciatura em Engenharia Informática\par}

    \vspace{1cm}

    {\Large Orientador:\par}
    \vspace{0.2cm}
    {\Large\bfseries Professor João Alfredo F. F. Dias, UBI\par}

    \vfill

    {\large \today \par}

\end{titlepage}

\blankpage

\pagenumbering{roman}

\chapter*{Agradecimentos}
\markboth{AGRADECIMENTOS}{}
\addcontentsline{toc}{chapter}{Agradecimentos}

Gostaria de agradecer ao meu orientador, Professor João Alfredo F. F. Dias, pelo acompanhamento, disponibilidade e orientação ao longo do desenvolvimento deste projeto. Agradeço também à Universidade da Beira Interior e ao Departamento de Informática pelas condições proporcionadas para a realização deste trabalho.

Um agradecimento especial à minha família e amigos pelo apoio constante durante esta etapa. O desenvolvimento deste projeto representou uma oportunidade de aplicar conhecimentos de computação gráfica, realidade estendida, jogos e engenharia de software a um contexto com potencial impacto social e humano.

\blankpage

\chapter*{Resumo}
\markboth{RESUMO}{}
\addcontentsline{toc}{chapter}{Resumo}

A doença de Alzheimer, nas suas fases iniciais, afeta progressivamente funções cognitivas como a memória, a atenção e a orientação espacial. Embora não exista ainda uma cura definitiva, intervenções não farmacológicas de estimulação cognitiva podem contribuir para manter capacidades funcionais, melhorar o envolvimento do utilizador e apoiar abordagens terapêuticas complementares. Neste contexto, a realidade virtual surge como uma tecnologia promissora, permitindo criar ambientes imersivos, controlados e seguros, adaptáveis ao ritmo do utilizador.

O presente projeto propõe o desenvolvimento de \textit{MindBridgeXR}, uma experiência em realidade estendida orientada para a estimulação cognitiva em fases iniciais da doença de Alzheimer. A aplicação decorre num ambiente doméstico virtual, onde o utilizador progride por diferentes fases: exploração inicial da casa, navegação guiada entre divisões, realização de um jogo de memória na sala de jantar e execução de uma tarefa funcional no exterior envolvendo seleção e entrega de alimentos. Estas tarefas foram concebidas para estimular memória visual, atenção, reconhecimento espacial, planeamento simples e execução de atividades próximas do quotidiano.

A solução foi implementada em Unity, recorrendo ao XR Interaction Toolkit e a uma arquitetura organizada em controladores de fases, zonas de divisão, transições entre cenas, objetos interativos e recolha automática de métricas. O projeto privilegia uma experiência acessível, com instruções curtas, feedback positivo e ausência de penalizações explícitas, procurando reduzir frustração e carga cognitiva excessiva. Este relatório descreve o enquadramento teórico, a conceção da solução, a implementação técnica e o procedimento de avaliação exploratória com utilizadores.

\textbf{Palavras-chave:} Realidade Virtual, Realidade Estendida, Alzheimer, Estimulação Cognitiva, Jogos Sérios, Unity.

\blankpage

\chapter*{Abstract}
\markboth{ABSTRACT}{}
\addcontentsline{toc}{chapter}{Abstract}

Alzheimer's disease, particularly in its early stages, progressively affects cognitive functions such as memory, attention, and spatial orientation. Although there is currently no definitive cure, non-pharmacological cognitive stimulation interventions may help preserve functional abilities, increase user engagement, and complement traditional therapeutic approaches. In this context, virtual reality is a promising technology, as it enables immersive, controlled, and safe environments that can be adapted to the user's pace.

This project presents \textit{MindBridgeXR}, an extended reality experience aimed at cognitive stimulation in the early stages of Alzheimer's disease. The application takes place in a virtual domestic environment where the user progresses through several phases: initial exploration of the house, guided navigation between rooms, a memory card game in the dining room, and an outdoor functional task involving the selection and delivery of food items. These tasks are designed to stimulate visual memory, attention, spatial recognition, simple planning, and execution of everyday-like activities.

The solution was implemented in Unity using the XR Interaction Toolkit and an architecture based on phase controllers, room zones, scene transitions, interactive objects, and automatic metrics collection. The project prioritizes accessibility through short instructions, positive feedback, and the absence of explicit failure states, aiming to reduce frustration and excessive cognitive load. This report describes the theoretical background, solution design, technical implementation, and the exploratory user evaluation procedure.

\textbf{Keywords:} Virtual Reality, Extended Reality, Alzheimer's Disease, Cognitive Stimulation, Serious Games, Unity.

\blankpage

\markboth{CONTEÚDO}{}
\tableofcontents

\blankpage

\markboth{LISTA DE FIGURAS}{}
\listoffigures

\blankpage

\markboth{LISTA DE TABELAS}{}
\listoftables

\blankpage

\chapter*{Acrónimos}
\markboth{ACRÓNIMOS}{}
\addcontentsline{toc}{chapter}{Acrónimos}
\begin{acronym}[XRI]
    \acro{AD}{Alzheimer's Disease}
    \acro{HMD}{Head-Mounted Display}
    \acro{TMP}{TextMesh Pro}
    \acro{UI}{User Interface}
    \acro{VR}{Virtual Reality}
    \acro{XR}{Extended Reality}
    \acro{XRI}{XR Interaction Toolkit}
\end{acronym}

\blankpage

\chapter*{Declaração de Uso de IA Generativa}
\markboth{DECLARAÇÃO DE USO DE IA GENERATIVA}{}
\addcontentsline{toc}{chapter}{Declaração de Uso de IA Generativa}

Durante a preparação deste relatório, foram utilizadas ferramentas de inteligência artificial generativa como apoio à estruturação textual, organização de secções, revisão linguística e formatação em \LaTeX. Todo o conteúdo foi revisto, adaptado e validado pelo autor, que assume total responsabilidade pela versão final do documento.

\blankpage

\clearpage
\pagenumbering{arabic}

\chapter{Introdução}

A doença de Alzheimer (\ac{AD}) é uma das principais causas de demência e representa um desafio crescente para os sistemas de saúde, para os cuidadores e para a qualidade de vida das pessoas afetadas. Nas fases iniciais, alterações na memória, atenção, orientação espacial e capacidade de executar atividades do quotidiano podem comprometer progressivamente a autonomia do indivíduo \cite{vasquez2025,oliveira2021}.

Apesar dos avanços médicos, as abordagens terapêuticas continuam a recorrer frequentemente a intervenções não farmacológicas, nomeadamente treino cognitivo, estimulação cognitiva, reminiscência e atividades funcionais adaptadas. A investigação recente tem destacado o potencial da \ac{VR} e da \ac{XR} para criar ambientes seguros, motivadores e ajustáveis, nos quais o utilizador pode realizar tarefas cognitivas de forma imersiva e controlada \cite{vasquez2025,goncalves2025}.

O projeto \textit{MindBridgeXR} enquadra-se neste domínio, propondo uma experiência imersiva num ambiente familiar, semelhante a uma casa, onde o utilizador realiza tarefas progressivas de exploração, orientação, memória e execução funcional.

\section{Motivação}

A principal motivação deste projeto é explorar o potencial da \ac{XR} como ferramenta de apoio à estimulação cognitiva em pessoas com \ac{AD} em fase inicial. Ambientes virtuais permitem simular situações do quotidiano com menor risco, maior controlo e possibilidade de adaptação. Ao contrário de testes cognitivos tradicionais, muitas vezes abstratos e pouco ecológicos, uma experiência virtual pode aproximar-se de tarefas reais, como localizar divisões, reconhecer objetos ou preparar uma mesa.

Além disso, a componente lúdica pode aumentar o envolvimento e reduzir a resistência à realização de exercícios cognitivos. A literatura analisada aponta que experiências imersivas podem contribuir para melhorias no envolvimento e, em alguns casos, no desempenho cognitivo e físico, embora a evidência disponível ainda seja heterogénea \cite{vasquez2025,yun2020,thapa2020}.

\section{Contextualização e Problema}

Muitos instrumentos de avaliação cognitiva decorrem em contextos controlados, com tarefas padronizadas que nem sempre representam os desafios do quotidiano. Estudos sobre diagnóstico gamificado em realidade virtual defendem que cenários familiares podem oferecer maior validade ecológica ao permitir observar comportamento em tarefas próximas da vida real \cite{xiong2024}.

No entanto, soluções em \ac{VR} para populações idosas e pessoas com défice cognitivo enfrentam desafios específicos: conforto visual, enjoo, complexidade dos controladores, excesso de estímulos, frustração perante erros e dificuldade em compreender instruções longas \cite{tuena2020,yun2020}. Assim, o problema central deste projeto consiste em conceber e implementar uma experiência \ac{XR} que seja cognitivamente estimulante, mas simultaneamente simples, confortável e emocionalmente segura.

\section{Objetivos}

O objetivo geral deste projeto é desenvolver um protótipo funcional em \ac{XR} para estimulação cognitiva em Alzheimer em fase inicial, usando tarefas curtas, progressivas e integradas num ambiente doméstico virtual.

Os objetivos específicos são:

\begin{itemize}
    \item Conceber um ambiente familiar e seguro, inspirado numa casa e no seu exterior.
    \item Implementar uma fase de exploração para familiarização espacial.
    \item Criar tarefas de navegação guiada entre divisões.
    \item Integrar um jogo de memória tridimensional na sala de jantar.
    \item Desenvolver uma tarefa funcional de preparação de alimentos no exterior.
    \item Usar feedback visual, textual e sonoro para orientar o utilizador.
    \item Evitar penalizações explícitas, privilegiando progressão e reforço positivo.
    \item Implementar a recolha automática de métricas de desempenho e interação ao longo das quatro fases.
\end{itemize}

\section{Contribuições}

As principais contribuições deste trabalho são:

\begin{itemize}
    \item Uma experiência \ac{XR} estruturada em fases cognitivas progressivas.
    \item Uma arquitetura em Unity baseada em controladores de fase e eventos.
    \item A integração de tarefas com maior validade ecológica, próximas de atividades de vida diária.
    \item Uma abordagem de design orientada para acessibilidade, conforto e segurança emocional.
    \item Um sistema de recolha e exportação de métricas por sessão, preparado para apoiar a avaliação com utilizadores.
\end{itemize}

\section{Planeamento e Cronograma}

O projeto foi planeado em fases, começando pela revisão bibliográfica e familiarização com Unity e \ac{XR}, seguindo para a construção do ambiente, implementação das tarefas, refinamento da experiência e preparação do relatório.

\begin{table}[H]
\centering
\caption{Cronograma geral do projeto}
\begin{tabular}{p{3cm}p{10cm}}
\toprule
\textbf{Período} & \textbf{Atividades principais} \\
\midrule
Fevereiro & Estado da arte, definição do problema, testes iniciais em Unity/XR. \\
Março & Construção do ambiente, locomoção, zonas da casa e interação básica. \\
Abril & Implementação das fases cognitivas, jogo de memória e tarefa de alimentos. \\
Maio & Refinamento, feedback visual/sonoro, testes internos e escrita técnica. \\
Junho & Consolidação do relatório, preparação da demonstração e entrega. \\
\bottomrule
\end{tabular}
\end{table}

\section{Estrutura do Documento}

Este documento encontra-se organizado em seis capítulos. O Capítulo 1 apresenta a introdução, motivação, problema e objetivos. O Capítulo 2 descreve e compara trabalhos relacionados sobre Alzheimer, estimulação cognitiva, realidade virtual e jogos sérios. O Capítulo 3 apresenta os requisitos, as fases e os princípios de conceção da solução. O Capítulo 4 descreve a implementação técnica, a interação, a arquitetura e o sistema de métricas. O Capítulo 5 apresenta a metodologia de avaliação, os resultados dos testes e a respetiva discussão. Por fim, o Capítulo 6 apresenta as conclusões, limitações e trabalho futuro.

\blankpage

\chapter{Estado da Arte}

\section{Doença de Alzheimer e Estimulação Cognitiva}

A doença de Alzheimer é uma condição neurodegenerativa progressiva que afeta sobretudo a memória, mas também funções como atenção, linguagem, orientação e tomada de decisão. Nas fases iniciais, o declínio pode ainda permitir a participação em atividades terapêuticas, tornando relevante a intervenção precoce \cite{oliveira2021,goncalves2025}.

A estimulação cognitiva procura manter ou ativar capacidades cognitivas através de tarefas estruturadas, repetidas e adaptadas ao utilizador. Estudos com pessoas com défice cognitivo ligeiro ou demência ligeira têm explorado tarefas imersivas e atividades motor-cognitivas como formas de treino e estimulação, reportando resultados promissores, embora ainda dependentes do desenho e duração de cada intervenção \cite{oliveira2021,thapa2020,yun2020}.

\section{Realidade Virtual em Contextos Terapêuticos}

A \ac{VR} permite criar ambientes tridimensionais imersivos com elevado controlo experimental. No contexto da saúde, esta tecnologia tem sido aplicada em treino cognitivo, reabilitação e simulação de atividades funcionais \cite{vasquez2025,goncalves2025}.

Em Alzheimer e défice cognitivo ligeiro, estudos recentes apontam que intervenções com \ac{VR}, exergaming e tarefas motor-cognitivas podem beneficiar a função cognitiva, a função física e, em menor grau, a qualidade de vida \cite{vasquez2025,thapa2020}. No entanto, a evidência ainda apresenta limitações devido à heterogeneidade dos protocolos, diferentes durações de intervenção, diferentes instrumentos de avaliação e amostras reduzidas \cite{vasquez2025}.

\section{Jogos Sérios e Gamificação}

Jogos sérios são aplicações lúdicas desenvolvidas com uma finalidade que ultrapassa o entretenimento. No caso da estimulação cognitiva, a gamificação pode favorecer a motivação, a repetição e o envolvimento do utilizador. Tarefas como jogos de memória, procura de objetos, navegação e associação podem ser integradas em ambientes virtuais para estimular diferentes processos cognitivos \cite{xiong2024,goncalves2025}.

Um princípio relevante para utilizadores idosos ou com défice cognitivo é a redução de frustração e de complexidade desnecessária. A literatura de usabilidade recomenda instruções claras, interfaces simples e apoio adequado às dificuldades de interação \cite{tuena2020,yun2020}. Este princípio influenciou diretamente o desenho do \textit{MindBridgeXR}, que privilegia continuidade, pistas graduais e reforço positivo.

\section{Validade Ecológica em Ambientes Virtuais}

Um problema recorrente em testes neuropsicológicos tradicionais é a sua limitada relação com o funcionamento real do utilizador em atividades quotidianas. Cenários virtuais como supermercados ou casas permitem observar comportamentos em tarefas mais naturais, mantendo controlo sobre estímulos, dificuldade e métricas \cite{xiong2024}.

O trabalho de Xiong et al.\ sobre diagnóstico gamificado em \ac{VR}, baseado numa tarefa de compras, demonstra esta direção: o utilizador memoriza uma lista, planeia uma ordem de compra e lida com distrações. A partir destas ações, podem ser recolhidas métricas relacionadas com memória verbal, memória de trabalho e funções executivas \cite{xiong2024}.

\section{Limitações e Desafios de Usabilidade}

Apesar do potencial da \ac{VR}, existem desafios importantes para utilizadores idosos: desconforto visual, dificuldade com controladores físicos, fadiga, risco de enjoo, excesso de estímulos e necessidade de instruções simples \cite{tuena2020}. A acessibilidade deve, por isso, ser considerada desde o início.

No \textit{MindBridgeXR}, estes desafios traduzem-se em decisões como tarefas curtas, ambiente familiar, feedback positivo, progressão faseada, uso de objetos reconhecíveis e redução da complexidade da interface.

\section{Síntese Comparativa dos Trabalhos Relacionados}

Os trabalhos analisados diferem quanto ao objetivo, população e desenho experimental. Alguns avaliam intervenções cognitivas ou motor-cognitivas, enquanto outros se concentram na viabilidade, usabilidade ou utilização da \ac{VR} como instrumento de avaliação. A Tabela~\ref{tab:estado-arte-comparacao} resume estas diferenças e identifica a influência de cada trabalho no desenvolvimento do \textit{MindBridgeXR}.

\begin{longtable}{p{2.2cm}p{2.8cm}p{3.2cm}p{3.7cm}}
\caption{Comparação dos principais trabalhos relacionados}
\label{tab:estado-arte-comparacao}\\
\toprule
\textbf{Trabalho} & \textbf{Tipo e população} & \textbf{Abordagem} & \textbf{Contributo para o MindBridgeXR} \\
\midrule
\endfirsthead
\toprule
\textbf{Trabalho} & \textbf{Tipo e população} & \textbf{Abordagem} & \textbf{Contributo para o MindBridgeXR} \\
\midrule
\endhead
Vásquez-Carrasco et al.\ \cite{vasquez2025} &
Revisão sistemática sobre pessoas idosas com doença de Alzheimer. &
Síntese de intervenções em \ac{VR} dirigidas à função cognitiva, função física e qualidade de vida. &
Sustenta o potencial da tecnologia, mas evidencia heterogeneidade metodológica e necessidade de prudência na interpretação dos resultados. \\

Xiong et al.\ \cite{xiong2024} &
Proposta de diagnóstico gamificado em \ac{VR}. &
Tarefa de compras com memorização de lista, planeamento do percurso e distrações. &
Inspirou a utilização de tarefas quotidianas, listas de objetos e métricas comportamentais com maior validade ecológica. \\

Gonçalves et al.\ \cite{goncalves2025} &
Experiência imersiva dirigida a fases iniciais de Alzheimer. &
Atividades de estimulação cognitiva integradas num ambiente virtual. &
Reforça a pertinência de tarefas imersivas progressivas adaptadas às capacidades do utilizador. \\

Oliveira et al.\ \cite{oliveira2021} &
Ensaio piloto aleatorizado com pessoas com demência ligeira a moderada devido a Alzheimer. &
Programa de estimulação cognitiva baseado em \ac{VR}. &
Demonstra a possibilidade de estruturar intervenções em sessões e combinar desempenho com avaliação cognitiva e funcional. \\

Yun et al.\ \cite{yun2020} &
Estudo de viabilidade e usabilidade com défice cognitivo ligeiro e demência ligeira. &
Treino cognitivo totalmente imersivo num ambiente enriquecido. &
Destaca a importância da familiarização, simplicidade de interação, conforto e observação de dificuldades de utilização. \\

Thapa et al.\ \cite{thapa2020} &
Ensaio controlado aleatorizado com adultos mais velhos com défice cognitivo ligeiro. &
Intervenção em \ac{VR} com componentes cognitivas e motoras. &
Mostra o interesse de combinar diferentes domínios de desempenho e registar medidas objetivas ao longo das tarefas. \\

Tuena et al.\ \cite{tuena2020} &
Revisão sistemática de aplicações clínicas e de investigação com pessoas idosas. &
Análise de problemas de usabilidade, tolerabilidade e interação em \ac{VR}. &
Fundamenta decisões de acessibilidade, instruções curtas, apoio do avaliador e monitorização de desconforto. \\
\bottomrule
\end{longtable}

Apesar dos resultados promissores, os estudos apresentam protocolos, equipamentos, durações e instrumentos de avaliação distintos. Esta diversidade dificulta comparações diretas e impede assumir que uma melhoria observada num contexto se transfere automaticamente para outro. Por esse motivo, o \textit{MindBridgeXR} é apresentado como um protótipo de apoio à estimulação cognitiva e não como uma ferramenta clínica validada.

\blankpage

\chapter{Conceção da Solução}

\section{Visão Geral}

O \textit{MindBridgeXR} é uma experiência imersiva em \ac{XR} que decorre num ambiente doméstico virtual. A aplicação organiza-se em quatro fases sequenciais, nas quais o utilizador passa de uma exploração livre para tarefas cada vez mais dirigidas.

\begin{figure}[H]
    \centering
   \includegraphics[width=\textwidth]{Fluxo_MindBridgeXR.png}
    \caption{Fluxo geral das fases do MindBridgeXR}
    \label{fig:fluxo-geral}
\end{figure}

\section{Requisitos da Solução}

Os requisitos foram definidos a partir dos objetivos do projeto, das limitações de usabilidade identificadas na literatura e das características do equipamento utilizado.

\subsection{Requisitos Funcionais}

\begin{itemize}
    \item Permitir a exploração das divisões da casa e registar as zonas visitadas.
    \item Apresentar tarefas de navegação guiada e validar a chegada à divisão solicitada.
    \item Disponibilizar um jogo de memória tridimensional com seleção de cartas e validação de pares.
    \item Apresentar uma lista de alimentos e validar individualmente as entregas realizadas.
    \item Coordenar automaticamente a progressão entre as quatro fases.
    \item Fornecer instruções e feedback visual, textual e sonoro durante a experiência.
    \item Guardar métricas de desempenho e interação associadas a uma sessão anónima.
    \item Preservar o estado necessário quando o utilizador muda de cena.
\end{itemize}

\subsection{Requisitos Não Funcionais}

\begin{itemize}
    \item \textbf{Usabilidade:} as ações principais devem exigir poucos comandos e ser demonstráveis antes do teste.
    \item \textbf{Acessibilidade:} os textos devem ser curtos, legíveis e acompanhados por pistas visuais sempre que necessário.
    \item \textbf{Conforto:} a velocidade de deslocação deve ser moderada e as mudanças de cena devem usar transições graduais.
    \item \textbf{Segurança emocional:} a aplicação deve evitar mensagens de derrota, penalizações agressivas e pressão temporal explícita.
    \item \textbf{Privacidade:} os dados devem ser associados a identificadores anónimos e guardados localmente.
    \item \textbf{Manutenibilidade:} a lógica das fases, interação e métricas deve permanecer separada em componentes com responsabilidades distintas.
    \item \textbf{Desempenho:} a experiência deve manter uma apresentação visual estável no dispositivo-alvo, evitando interrupções percetíveis durante a interação.
\end{itemize}

\subsection{Restrições}

O protótipo foi concebido para o Meta Quest 3, utilizando os respetivos comandos e uma área física livre de obstáculos. A experiência depende da capacidade do participante para compreender instruções simples, movimentar os membros superiores e tolerar a utilização de um \ac{HMD}. O espaço virtual encontra-se distribuído por três cenas, o que exige transições e reposicionamento controlado do utilizador.

\section{Fase 1: Exploração Tutorial}

Na primeira fase, o utilizador explora livremente a casa. Sempre que entra numa divisão, o sistema identifica a zona, apresenta temporariamente o nome da divisão e marca-a como visitada. Esta fase serve para familiarizar o utilizador com o espaço antes de exigir tarefas específicas.

\section{Fase 2: Navegação Guiada}

Após visitar todas as divisões, o sistema inicia uma fase de navegação guiada. O utilizador recebe instruções como ``vai para a sala de jantar'' ou ``vai para a cozinha''. As tarefas são selecionadas a partir de uma lista global de divisões, podendo ser baralhadas e limitadas por ronda.

Esta fase estimula orientação espacial, compreensão de instruções e memória de localização.

\section{Fase 3: Jogo de Memória na Sala de Jantar}

Na terceira fase, o utilizador é instruído a dirigir-se à sala de jantar e aproximar-se da mesa. Ao chegar, inicia-se um jogo de memória com cartas tridimensionais. O jogo apresenta cartas com pares associados e o utilizador deve selecionar duas cartas de cada vez, procurando correspondências.

Esta tarefa estimula memória visual, atenção e reconhecimento de padrões. Antes da ronda começar, as cartas podem ser reveladas por um curto período, funcionando como fase de memorização inicial.

\section{Fase 4: Preparação de Alimentos no Exterior}

Na fase final, o utilizador vai para o exterior, pega numa lista de alimentos e deve colocar na mesa os alimentos pedidos. O sistema valida cada entrega, aceita alimentos corretos e rejeita alimentos que não pertencem à lista ou excedem a quantidade pedida.

Esta tarefa aproxima-se de uma atividade funcional do quotidiano, envolvendo leitura ou interpretação de uma lista, seleção de objetos, memória de objetivo e execução motora simples.

\section{Relação entre Tarefas, Funções Cognitivas e Métricas}

As tarefas foram concebidas para mobilizar diferentes processos cognitivos e funcionais. Esta relação representa uma intenção de design fundamentada na literatura, não uma demonstração de eficácia clínica. A avaliação permite observar o comportamento durante as tarefas, mas não substitui instrumentos neuropsicológicos validados.

\begin{table}[H]
\centering
\caption{Relação entre tarefas, processos mobilizados e métricas}
\label{tab:tarefas-cognicao-metricas}
\begin{tabular}{p{2.1cm}p{2.8cm}p{3.2cm}p{3.7cm}}
\toprule
\textbf{Fase} & \textbf{Ação principal} & \textbf{Processos mobilizados} & \textbf{Métricas associadas} \\
\midrule
Exploração & Descobrir e visitar as divisões. & Orientação espacial, atenção ao ambiente e aprendizagem do espaço. & Ordem de visita, divisões únicas, revisitas, tempo e distância percorrida. \\
Navegação guiada & Encontrar uma divisão indicada. & Memória espacial, compreensão de instruções e tomada de decisão. & Tempo por tarefa, percurso, mudanças de cena e regressos. \\
Jogo de memória & Selecionar cartas e encontrar pares. & Memória visual, atenção e reconhecimento de padrões. & Tentativas, acertos, erros, duração, taxa de acerto e eficiência. \\
Alimentos & Interpretar a lista, selecionar e entregar objetos. & Atenção, memória de objetivo, planeamento simples e execução motora. & Objetos agarrados, entregas, rejeições, tempo, precisão e manipulações desnecessárias. \\
\bottomrule
\end{tabular}
\end{table}

\section{Princípios de Design}

A solução foi orientada pelos seguintes princípios:

\begin{itemize}
    \item \textbf{Ambiente familiar:} uso de uma casa e exterior como contexto reconhecível.
    \item \textbf{Progressão gradual:} início com exploração livre e aumento gradual da exigência.
    \item \textbf{Instruções simples:} mensagens curtas e objetivas.
    \item \textbf{Feedback positivo:} reforço de progresso sem penalização agressiva.
    \item \textbf{Segurança emocional:} ausência de ecrãs de derrota ou falha explícita.
    \item \textbf{Validade ecológica:} tarefas inspiradas em ações reais do quotidiano.
\end{itemize}

\blankpage

\chapter{Implementação}

\section{Tecnologias Utilizadas}

O projeto foi desenvolvido em Unity \cite{unity}, com recurso ao \ac{XRI} para suportar interação em \ac{VR}. A implementação foi organizada em scripts C\#, responsáveis pela gestão das fases, deteção de divisões, transições entre cenas e interação com objetos.

\begin{table}[H]
\centering
\caption{Tecnologias principais}
\begin{tabular}{ll}
\toprule
\textbf{Tecnologia} & \textbf{Utilização} \\
\midrule
Unity & Motor de jogo e ambiente de desenvolvimento 3D. \\
C\# & Linguagem usada nos scripts de comportamento. \\
\ac{XRI} & Interação, seleção e manipulação de objetos em \ac{XR}. \\
\ac{TMP} & Apresentação de instruções e mensagens ao utilizador. \\
Meta Quest / HMD & Plataforma-alvo para execução da experiência. \\
\bottomrule
\end{tabular}
\end{table}

\section{Configuração Técnica}

A versão utilizada durante o desenvolvimento foi Unity 6.3, concretamente \texttt{6000.3.7f1}. O projeto utiliza o \ac{XRI} 3.3.1, Input System 1.18.0, OpenXR 1.16.1, Meta OpenXR 2.3.1 e Universal Render Pipeline 17.3.0. O destino principal de compilação é Android com arquitetura ARM64, compatível com execução autónoma no Meta Quest 3.

\begin{table}[H]
\centering
\caption{Configuração técnica principal do protótipo}
\label{tab:configuracao-tecnica}
\begin{tabular}{p{5cm}p{8cm}}
\toprule
\textbf{Componente} & \textbf{Configuração} \\
\midrule
Motor & Unity 6.3 (\texttt{6000.3.7f1}) \\
Renderização & Universal Render Pipeline 17.3.0 \\
Interação XR & XR Interaction Toolkit 3.3.1 \\
Entrada & Unity Input System 1.18.0 \\
Runtime XR & OpenXR 1.16.1 e Meta OpenXR 2.3.1 \\
Dispositivo-alvo & Meta Quest 3 com comandos Touch Plus \\
Plataforma & Android, arquitetura ARM64 \\
Cenas incluídas & \texttt{Floor1\_Scene}, \texttt{Floor2\_Scene} e \texttt{Exterior\_Scene} \\
\bottomrule
\end{tabular}
\end{table}

O perfil de qualidade móvel reduz a complexidade de alguns efeitos relativamente ao perfil de computador. Não foi implementado um limite manual de fotogramas por segundo, ficando a sincronização dependente do runtime XR e da frequência de atualização ativa no dispositivo. A avaliação de desempenho deverá, por isso, confirmar a estabilidade visual diretamente no Meta Quest 3 e registar eventuais quebras, atrasos ou desconforto associados.

\section{Interação e Locomoção}

O utilizador desloca-se através de movimento contínuo controlado pelo joystick. O componente de locomoção foi configurado com uma velocidade de aproximadamente \(0{,}85\) unidades por segundo, inferior ao valor original do \textit{prefab}, procurando tornar o movimento mais previsível. A orientação pode ser alterada fisicamente através da cabeça e do corpo ou por rotação em incrementos de \(45^\circ\), reduzindo a rotação virtual contínua.

A manipulação dos alimentos e da lista utiliza objetos do tipo \texttt{XRGrabInteractable}. O utilizador aproxima a mão virtual e mantém pressionado o botão lateral do comando para agarrar, libertando-o para largar. As cartas do jogo de memória utilizam \texttt{XRSimpleInteractable} e são ativadas por seleção, evitando exigir a sua deslocação física.

As pistas de orientação combinam diferentes modalidades:

\begin{itemize}
    \item texto com a tarefa atual;
    \item destaque visual da divisão solicitada quando esta se encontra na cena;
    \item seta tridimensional sobre a lista de alimentos;
    \item destaque da mesa após a recolha da lista;
    \item mensagens centrais e sons de confirmação ou rejeição;
    \item transições com \textit{fade} durante mudanças de cena e reposicionamentos.
\end{itemize}

Esta combinação procura evitar que o sucesso dependa exclusivamente de texto, som ou memória das instruções. A interação mantém-se, contudo, dependente de controladores físicos, constituindo uma possível dificuldade para participantes sem experiência prévia.

\section{Recursos Visuais e Assets}

Para acelerar o desenvolvimento do protótipo e permitir maior foco na lógica da experiência, foram utilizados recursos visuais provenientes de repositórios e lojas de assets. O ambiente doméstico base foi apoiado por um modelo modular de casa disponível no CGTrader \cite{cgtraderhouse}. Alguns objetos específicos foram obtidos através do Sketchfab, incluindo um indicador em forma de seta \cite{sketchfabarrow}, um conjunto de peças para jogo de memória \cite{sketchfabmemory}, um cabide \cite{sketchfabcoathanger} e um comando de televisão \cite{sketchfabremote}. Para os alimentos usados na tarefa exterior, foi utilizado o pacote \textit{Farm Food Pack} da Unity Asset Store \cite{unityfarmfood}. Adicionalmente, foram usados recursos sonoros provenientes do Pixabay e da plataforma MyInstants, aplicados a impactos, transições e feedback de resposta correta ou incorreta \cite{soundswooshsuave,soundthudheavy,soundvaseimpact,soundpanhit,soundcorrect,soundwrong,soundduolingocorrect,soundduolingowrong}. A música ambiente utilizada foi \textit{The Autumn Wind}, de Wayne Lundh Trio \cite{theautumnwind}.

Estes recursos foram integrados no Unity e adaptados ao cenário, mantendo o objetivo de criar um ambiente familiar, reconhecível e visualmente simples para o utilizador.

\section{Arquitetura Geral}

A arquitetura da aplicação combina um coordenador global com componentes especializados e scripts auxiliares de responsabilidade limitada. O \texttt{TaskManager} conhece a sequência completa da experiência, mas delega o funcionamento interno do jogo de memória, da tarefa de alimentos, das transições e das métricas nos respetivos módulos. Esta separação reduz a dimensão de cada responsabilidade e permite alterar uma tarefa sem reimplementar toda a progressão.

\subsection{Coordenação, Divisões e Cenas}

\begin{longtable}{p{4.2cm}p{8.5cm}}
\caption{Scripts de coordenação, divisões e cenas}
\label{tab:scripts-coordenacao}\\
\toprule
\textbf{Script} & \textbf{Responsabilidade} \\
\midrule
\endfirsthead
\toprule
\textbf{Script} & \textbf{Responsabilidade} \\
\midrule
\endhead
\texttt{TaskManager} &
Mantém a fase atual, conhece o catálogo global de divisões, coordena a progressão entre as quatro fases, apresenta instruções e ativa os controladores especializados. Inicia também a sessão de métricas e comunica os eventos relevantes ao \texttt{MetricsManager}. \\

\texttt{RoomZone} &
Representa uma divisão através de um identificador e de uma zona de colisão. Notifica o \texttt{TaskManager} quando o jogador entra, controla o destaque do destino e apresenta a marca de divisão já visitada. \\

\texttt{SceneRoomRegistry} &
Regista no \texttt{TaskManager} as instâncias de \texttt{RoomZone} existentes na cena atual, permitindo atualizar destaques e marcas após cada carregamento. \\

\texttt{SceneTransitionTrigger} &
Deteta a entrada do jogador numa passagem entre espaços e solicita ao \texttt{SceneTransitionManager} o carregamento da cena e do ponto de destino correspondentes. \\

\texttt{SceneTransitionManager} &
Executa o \textit{fade}, apresenta mensagens de transição, carrega cenas e reposiciona o \texttt{XROrigin} no ponto de entrada. Reativa ainda a locomoção e os controladores após o carregamento. \\

\texttt{SceneSpawnPoint} &
Identifica um ponto de entrada através de um \texttt{spawnID}, usado pelo gestor de transições para determinar a posição e orientação do jogador. \\

\texttt{PersistentXRRig} &
Impede a duplicação do rig \ac{XR} e permite que a instância principal permaneça ativa durante mudanças de cena. \\
\bottomrule
\end{longtable}

\subsection{Jogo de Memória}

\begin{longtable}{p{4.2cm}p{8.5cm}}
\caption{Scripts do jogo de memória}
\label{tab:scripts-memoria}\\
\toprule
\textbf{Script} & \textbf{Responsabilidade} \\
\midrule
\endfirsthead
\toprule
\textbf{Script} & \textbf{Responsabilidade} \\
\midrule
\endhead
\texttt{DiningMemoryPhaseController} &
Implementa a máquina de estados da terceira fase: chegada à sala de jantar, aproximação à mesa, execução do jogo e conclusão. Liga os eventos espaciais ao controlador do minijogo e às métricas. \\

\texttt{DiningTableZone} &
Deteta a aproximação do jogador à mesa e emite uma única notificação, que permite iniciar o jogo no momento correto. \\

\texttt{MemoryMiniGame3DController} &
Descobre e prepara as cartas, executa a pré-visualização inicial, recebe seleções, compara os pares, controla os atrasos entre tentativas e determina a conclusão da ronda. \\

\texttt{MemoryCard3D} &
Guarda o identificador do par e o estado individual de uma carta. Controla a animação de revelação e ocultação, impede novas seleções quando necessário e emite o evento de seleção. \\

\texttt{XRMemoryCardSelectBridge} &
Faz a ponte entre o evento de seleção do \texttt{XRSimpleInteractable} e a lógica da \texttt{MemoryCard3D}, mantendo a dependência do XRI fora do comportamento principal da carta. \\
\bottomrule
\end{longtable}

\subsection{Tarefa de Alimentos}

\begin{longtable}{p{4.2cm}p{8.5cm}}
\caption{Scripts e tipos da tarefa de alimentos}
\label{tab:scripts-alimentos}\\
\toprule
\textbf{Script ou tipo} & \textbf{Responsabilidade} \\
\midrule
\endfirsthead
\toprule
\textbf{Script ou tipo} & \textbf{Responsabilidade} \\
\midrule
\endhead
\texttt{OutdoorFoodPhaseController} &
Implementa a máquina de estados da fase exterior, mantém a lista e as quantidades pedidas, valida cada tentativa de entrega, atualiza os pratos, apresenta pistas e emite a conclusão da fase. \\

\texttt{FoodCollectible} &
Representa um alimento manipulável. Guarda o tipo e o identificador usado nas métricas, acompanha os eventos de agarrar e largar, preserva a posição inicial e executa o retorno de objetos rejeitados. \\

\texttt{FoodListPickup} &
Deteta a primeira recolha da lista, remove a seta de indicação e notifica o controlador para iniciar a etapa de entrega. \\

\texttt{FoodDeliveryZone} &
Deteta alimentos colocados na mesa, espera que sejam largados, evita processamentos duplicados e aplica o resultado devolvido pelo controlador da fase. \\

\texttt{WorldArrowIndicator} &
Orienta visualmente o utilizador para um alvo, rodando na sua direção e aplicando um movimento vertical periódico para aumentar a visibilidade. \\

\texttt{FoodType} &
Enumeração que identifica as categorias de alimentos reconhecidas pelo sistema, permitindo comparar os objetos entregues com os requisitos da lista. \\

\texttt{FoodDeliveryResult} &
Enumeração que representa o resultado da validação de uma entrega: ignorada, aceite ou rejeitada com retorno à posição inicial. \\
\bottomrule
\end{longtable}

\subsection{Métricas e Avaliação}

\begin{longtable}{p{4.2cm}p{8.5cm}}
\caption{Scripts do sistema de métricas}
\label{tab:scripts-metricas}\\
\toprule
\textbf{Script} & \textbf{Responsabilidade} \\
\midrule
\endfirsthead
\toprule
\textbf{Script} & \textbf{Responsabilidade} \\
\midrule
\endhead
\texttt{ParticipantIdEntryUI} &
Cria uma interface tridimensional para introdução do identificador anónimo. Suspende temporariamente a locomoção, valida o código e inicia a sessão após confirmação. \\

\texttt{MetricsManager} &
Mantém os dados da sessão e das quatro fases, calcula tempos e métricas derivadas, deteta interrupções e mudanças de cena e guarda eventos e resumos em ficheiros persistentes. \\

\texttt{MetricsData} &
Define as classes serializáveis usadas para representar sessões, fases, visitas, tarefas guiadas, tentativas de memória, entregas de alimentos e eventos cronológicos. \\

\texttt{MetricsReportExporter} &
Lê os resumos das sessões e produz ficheiros comparativos em CSV e TSV, além de relatórios textuais detalhados para cada participante. \\
\bottomrule
\end{longtable}

\subsection{Interface e Áudio}

\begin{longtable}{p{4.2cm}p{8.5cm}}
\caption{Scripts auxiliares de interface e áudio}
\label{tab:scripts-auxiliares}\\
\toprule
\textbf{Script} & \textbf{Responsabilidade} \\
\midrule
\endfirsthead
\toprule
\textbf{Script} & \textbf{Responsabilidade} \\
\midrule
\endhead
\texttt{BackgroundMusicManager} &
Mantém a música ambiente entre cenas, configura a reprodução em ciclo e permite controlar o volume global. \\

\texttt{GrabSound} &
Reproduz um efeito sonoro quando um objeto com \texttt{XRGrabInteractable} é agarrado. \\

\texttt{ImpactSound} &
Reproduz sons de impacto de acordo com a velocidade relativa da colisão, aplicando um limiar mínimo e um intervalo entre reproduções. \\

\texttt{LookAtPlayer} &
Mantém elementos informativos orientados horizontalmente para a câmara do jogador, facilitando a leitura no espaço tridimensional. \\
\bottomrule
\end{longtable}

\subsection{Fluxo de Comunicação}

O \texttt{TaskManager} funciona como coordenador principal. As instâncias de \texttt{RoomZone} comunicam entradas em divisões, permitindo atualizar a fase atual e as métricas. O \texttt{SceneRoomRegistry} renova as referências locais após cada carregamento, enquanto os pares \texttt{SceneTransitionTrigger}--\texttt{SceneTransitionManager} e \texttt{SceneSpawnPoint} realizam a deslocação entre cenas.

Quando uma fase especializada começa, o gestor ativa o controlador correspondente e subscreve o respetivo evento de conclusão. Na fase de memória, a sequência principal é \texttt{DiningTableZone}, \texttt{DiningMemoryPhaseController}, \texttt{MemoryMiniGame3DController} e \texttt{MemoryCard3D}. Na fase exterior, os eventos de \texttt{FoodListPickup}, \texttt{FoodCollectible} e \texttt{FoodDeliveryZone} convergem no \texttt{OutdoorFoodPhaseController}.

O \texttt{MetricsManager} permanece ativo entre cenas e recebe notificações dos vários módulos sem controlar diretamente a lógica da experiência. Esta separação reduz o acoplamento entre avaliação e progressão. Os dados seguem as estruturas definidas em \texttt{MetricsData} e são posteriormente transformados pelo \texttt{MetricsReportExporter}. Os scripts de áudio e orientação visual reagem localmente aos eventos de interação, sem interferirem com as regras de progressão.

\section{Gestão Global de Fases}

O \texttt{TaskManager} mantém o estado atual da experiência através de uma enumeração com quatro fases principais: exploração tutorial, navegação guiada, memória na sala de jantar e alimentos no exterior. Este gestor recebe notificações quando o jogador entra numa divisão e decide que ação tomar de acordo com o estado atual.

\begin{lstlisting}[language=CSharp,caption={Estados principais da experiência}]
private enum GamePhase
{
    TutorialExploration,
    GuidedNavigation,
    DiningMemory,
    OutdoorFood
}
\end{lstlisting}

Durante a fase de tutorial, o sistema acumula as divisões visitadas. Quando todas são visitadas, inicia a navegação guiada. Depois da conclusão das tarefas de navegação, o sistema inicia a fase de memória. Por fim, quando a ronda de memória termina, a fase de alimentos é ativada.

\section{Deteção de Divisões}

Cada divisão possui uma zona de colisão com o script \texttt{RoomZone}. Quando o jogador entra na zona, o script comunica com o \texttt{TaskManager}. Esta abordagem permite separar a lógica de deteção espacial da lógica de progressão da experiência.

\begin{lstlisting}[language=CSharp,caption={Comunicação de entrada numa divisão}]
private void OnTriggerEnter(Collider other)
{
    if (!other.CompareTag("Player"))
        return;

    if (TaskManager.Instance != null)
        TaskManager.Instance.PlayerEnteredRoom(this);
}
\end{lstlisting}

\section{Jogo de Memória}

O jogo de memória utiliza cartas tridimensionais interativas. O controlador guarda a primeira e a segunda carta selecionadas, compara os respetivos identificadores de par e decide se as cartas ficam concluídas ou se voltam a ficar ocultas.

A fase inclui uma breve pré-visualização inicial das cartas, permitindo ao utilizador memorizar as posições antes do início da interação. Esta decisão reduz a frustração e torna a tarefa mais adequada a um contexto de estimulação.

\section{Tarefa de Alimentos}

A fase exterior é implementada como uma máquina de estados com quatro momentos: deslocação até ao exterior, procura da lista, entrega de alimentos e conclusão. Antes de pegar na lista, o utilizador é orientado a fazê-lo; depois da lista ser recolhida, a mesa passa a ser destacada como local de entrega.

Cada alimento possui um tipo, uma posição inicial e um comportamento de retorno. Se o alimento entregue for inválido, regressa ao ponto inicial. Se for aceite, é marcado como entregue e removido da cena.

\section{Transições entre Cenas}

O \texttt{SceneTransitionManager} gere transições com fade, carregamento de cena e reposicionamento do rig \ac{XR} num ponto de spawn. Esta funcionalidade permite reiniciar a posição do utilizador entre fases, garantindo maior controlo sobre o percurso da experiência.

\section{Interface e Feedback}

A \ac{UI} usa textos simples para indicar a tarefa atual. Mensagens centrais temporárias reforçam eventos importantes, como entrada numa divisão, recolha da lista ou conclusão de uma tarefa. A solução prevê ainda feedback sonoro para aceitação, rejeição, conclusão, transições suaves e impactos de objetos. Estes sons foram obtidos em bibliotecas online e integrados no Unity de forma a tornar as interações mais percetíveis, sem sobrecarregar o utilizador.

\section{Recolha e Exportação de Métricas}

A aplicação integra um sistema de recolha automática de métricas através do componente \texttt{MetricsManager}. Este componente é criado no início da experiência, permanece ativo durante as mudanças de cena e associa os dados a um identificador anónimo de participante e a um identificador único de sessão. São também registados a versão da aplicação, a plataforma, o modelo do dispositivo, a duração total da experiência, as mudanças de cena, eventuais interrupções e a conclusão de cada fase.

A recolha é organizada de acordo com as quatro fases da experiência:

\begin{itemize}
    \item \textbf{Exploração tutorial:} sequência de divisões visitadas, tempo até à primeira divisão, número de visitas e revisitas, tempo passado em cada divisão, mudanças de cena e distância percorrida.
    \item \textbf{Navegação guiada:} destino de cada tarefa, sequência de divisões percorrida, duração, mudanças de cena e regressos ou revisitas antes de alcançar o destino.
    \item \textbf{Jogo de memória:} tempo até à sala e à mesa, duração do jogo, cartas selecionadas, número e duração das tentativas, pares corretos e incorretos, taxa de acerto e eficiência relativamente ao mínimo teórico de tentativas.
    \item \textbf{Tarefa de alimentos:} tempo até à recolha da lista, objetos agarrados, tentativas de entrega, entregas aceites e rejeitadas, motivos de rejeição, manipulações desnecessárias, taxa de acerto e tempo médio por alimento correto.
\end{itemize}

Os eventos são registados cronologicamente num ficheiro \texttt{JSONL}, enquanto o resumo de cada sessão é guardado em \texttt{JSON}. O sistema gera ainda ficheiros \texttt{CSV} e \texttt{TSV} para comparação entre participantes, bem como relatórios textuais individuais. Os ficheiros são armazenados na pasta persistente da aplicação, permitindo preservar os dados entre execuções e facilitar a sua análise posterior.

Para proteger a privacidade dos participantes, o sistema não necessita de guardar nomes ou outros dados diretamente identificáveis. A correspondência entre o identificador anónimo e o participante, caso seja necessária durante o estudo, deverá ser mantida separadamente e com acesso restrito.

\subsection{Definição das Métricas Derivadas}

No jogo de memória, a taxa de acerto corresponde à proporção de tentativas que resultaram num par correto:

\begin{equation}
\text{Taxa de acerto}_{\text{memória}} =
\frac{N_{\text{pares corretos}}}{N_{\text{tentativas}}}\times 100.
\end{equation}

A eficiência compara o número mínimo teórico de tentativas, equivalente ao número de pares, com o número efetivamente realizado:

\begin{equation}
\text{Eficiência}_{\text{memória}} =
\min\left(1,\frac{N_{\text{tentativas mínimas}}}{N_{\text{tentativas realizadas}}}\right)\times 100.
\end{equation}

Na fase de alimentos, a taxa de acerto é calculada através da proporção de entregas aceites:

\begin{equation}
\text{Taxa de acerto}_{\text{alimentos}} =
\frac{N_{\text{entregas aceites}}}{N_{\text{tentativas de entrega}}}\times 100.
\end{equation}

As manipulações desnecessárias correspondem à diferença, limitada a valores não negativos, entre o número total de vezes que alimentos foram agarrados e o número de entregas aceites. Na navegação guiada, uma revisita ou regresso é contabilizado quando o participante volta a entrar numa divisão já visitada durante a mesma tarefa. A distância percorrida é estimada através da soma das distâncias entre posições sucessivas do jogador, recolhidas a intervalos configuráveis; esta métrica pode ser desativada e deve ser interpretada como uma aproximação.

\section{Desafios de Implementação}

O desenvolvimento exigiu resolver problemas que ultrapassam a implementação isolada de cada tarefa:

\begin{itemize}
    \item \textbf{Persistência entre cenas:} o gestor de tarefas, o sistema de transições e as métricas têm de sobreviver ao carregamento das três cenas sem criar instâncias duplicadas.
    \item \textbf{Reposicionamento do rig XR:} após uma mudança de cena, a posição da câmara deve coincidir com o ponto de entrada, respeitando simultaneamente o deslocamento físico do headset dentro do \texttt{XROrigin}.
    \item \textbf{Reinicialização da interação:} os objetos de locomoção e os controladores são temporariamente desativados e reativados após uma transição para evitar estados inconsistentes.
    \item \textbf{Validação de alimentos:} a zona de entrega deve impedir múltiplos registos causados por vários colisores, controlar quantidades e devolver objetos incorretos à posição inicial.
    \item \textbf{Continuidade do progresso:} o estado da tarefa exterior deve ser reconstruído quando ocorre uma mudança de cena, incluindo quantidades entregues e representações nos pratos.
    \item \textbf{Sessões incompletas:} o sistema de métricas deteta interrupções, perda de foco e reinícios posteriores a sessões que não chegaram ao fim.
    \item \textbf{Exportação de dados:} os eventos detalhados e os resumos agregados têm formatos diferentes, sendo necessário produzir ficheiros comparáveis sem perder o histórico individual.
\end{itemize}

\blankpage

\chapter{Avaliação e Discussão}

\section{Plano de Avaliação}

A avaliação deverá começar com utilizadores saudáveis, por razões éticas e de segurança, antes de qualquer teste com pessoas diagnosticadas com \ac{AD}. Esta abordagem permite avaliar inicialmente a usabilidade, o conforto, a clareza das instruções e a robustez técnica do protótipo, aspetos destacados como relevantes em estudos de viabilidade e revisões de usabilidade de \ac{VR} com pessoas idosas \cite{yun2020,tuena2020}.

O estudo deverá combinar observação direta, perguntas finais, escalas de avaliação subjetiva e as métricas automáticas descritas no capítulo anterior. O procedimento proposto encontra-se no Apêndice A, incluindo a apresentação do equipamento, o treino inicial, as tarefas, as frases de apoio permitidas e as observações a registar. Antes de começar, cada participante deverá ser informado sobre o objetivo do teste, o caráter voluntário da participação e a possibilidade de interromper a experiência a qualquer momento.

\subsection{Participantes}

A primeira avaliação destina-se a uma amostra de conveniência composta por adultos sem diagnóstico conhecido de doença de Alzheimer. A caracterização final deverá indicar o número de participantes, idade ou intervalo etário, género quando pertinente, experiência anterior com \ac{VR} e familiaridade com videojogos. Estes fatores são relevantes porque podem influenciar a utilização dos comandos, a velocidade de adaptação e o desempenho nas tarefas.

Antes da sessão, deve confirmar-se que o participante compreendeu o procedimento, aceita participar voluntariamente e não apresenta desconforto no momento inicial. A sessão deve ser interrompida perante pedido do participante ou sinais de náusea, tontura, fadiga acentuada, desorientação ou ansiedade.

\subsection{Procedimento}

Cada sessão deve seguir uma ordem uniforme:

\begin{enumerate}
    \item apresentação do estudo e recolha do consentimento;
    \item atribuição de um identificador anónimo;
    \item breve caracterização do participante e da sua experiência prévia;
    \item ajuste do headset e explicação dos comandos;
    \item treino inicial de aproximadamente um minuto;
    \item realização sequencial das quatro fases;
    \item interrupção imediata caso exista desconforto relevante;
    \item remoção do equipamento, perguntas finais e escala de avaliação;
    \item confirmação e armazenamento dos ficheiros de métricas.
\end{enumerate}

O avaliador deve evitar antecipar soluções. Quando for necessário ajudar, deve utilizar as frases neutras definidas no guião e registar o momento e o tipo de apoio prestado. A duração total, o local, as condições do espaço e quaisquer desvios ao procedimento deverão ser descritos juntamente com os resultados.

\subsection{Estratégia de Análise}

As métricas quantitativas deverão ser resumidas através de medidas descritivas adequadas, como média, mediana, mínimo, máximo e dispersão. Com amostras pequenas, os valores individuais e a mediana podem ser mais informativos do que a média isolada. As respostas em escala podem ser apresentadas por afirmação, enquanto os comentários abertos e as observações devem ser agrupados por temas, como locomoção, compreensão das instruções, manipulação, conforto e preferência de tarefa.

A análise deve procurar convergência entre fontes. Por exemplo, um tempo elevado numa fase pode representar dificuldade, exploração deliberada ou problema de interação; a interpretação só deverá ser feita após consultar o percurso, os pedidos de ajuda e os comentários do participante. Não se pretende realizar inferência clínica nem comparar o desempenho com valores normativos.

\section{Métricas de Avaliação}

\begin{table}[H]
\centering
\caption{Métricas para avaliação da experiência}
\begin{tabular}{p{4cm}p{9cm}}
\toprule
\textbf{Dimensão} & \textbf{Métrica} \\
\midrule
Exploração e orientação & Divisões visitadas, revisitas, sequência do percurso, distância percorrida, tempo até chegar ao destino e pedidos de ajuda. \\
Memória visual & Número de tentativas, pares corretos e incorretos, taxa de acerto, eficiência e tempo para concluir a ronda. \\
Atenção e execução & Alimentos agarrados, entregas aceites e rejeitadas, manipulações desnecessárias, taxa de acerto e tempo até completar a lista. \\
Usabilidade & Facilidade percebida, conforto, clareza das instruções, vontade de repetir. \\
Conforto em XR & Náusea, fadiga visual, desorientação, dificuldade com interação. \\
\bottomrule
\end{tabular}
\end{table}

\section{Expectativas de Avaliação}

Espera-se que a aplicação seja compreensível e suficientemente simples para ser usada em sessões curtas. Espera-se também que o ambiente doméstico favoreça a familiaridade e que as tarefas sejam percebidas como mais naturais do que testes cognitivos tradicionais. Com base na literatura, antecipa-se que experiências imersivas adequadamente desenhadas possam favorecer o envolvimento e a disponibilidade para repetir exercícios \cite{goncalves2025,yun2020}. Estas expectativas deverão ser confrontadas com os resultados apresentados na secção seguinte, sem serem interpretadas como evidência de eficácia terapêutica.

\section{Resultados dos Testes}

% Secção reservada para apresentação e análise dos resultados obtidos nos testes com utilizadores.

\section{Discussão}

O \textit{MindBridgeXR} procura responder a vários pontos identificados no estado da arte. Em primeiro lugar, utiliza um ambiente familiar e tarefas próximas da vida diária, procurando aumentar a validade ecológica face a exercícios cognitivos abstratos \cite{xiong2024}. Em segundo lugar, combina diferentes domínios cognitivos, como orientação espacial, memória visual, atenção e execução de instruções. Em terceiro lugar, adota uma abordagem emocionalmente segura, evitando penalizações explícitas e privilegiando progressão e feedback positivo.

Ainda assim, existem limitações. O protótipo necessita de validação com utilizadores para confirmar conforto, clareza das instruções e adequação da dificuldade. A ausência de testes clínicos impede qualquer conclusão sobre eficácia terapêutica. Embora o sistema já recolha dados objetivos de interação e desempenho, estas métricas devem ser interpretadas em conjunto com observação e respostas subjetivas, não constituindo por si só medidas clínicas de função cognitiva.

\section{Limitações do Trabalho}

As principais limitações identificadas são:

\begin{itemize}
    \item a avaliação inicial com participantes saudáveis não representa diretamente as capacidades, necessidades ou respostas de pessoas com doença de Alzheimer;
    \item a utilização de uma amostra de conveniência reduz a possibilidade de generalização;
    \item experiência anterior com videojogos ou \ac{VR} pode influenciar o desempenho independentemente das funções cognitivas mobilizadas;
    \item a ordem fixa das fases introduz possíveis efeitos de aprendizagem, adaptação e fadiga;
    \item os controladores físicos podem aumentar a dificuldade motora e tecnológica;
    \item as métricas comportamentais não constituem instrumentos de diagnóstico nem medidas clínicas validadas;
    \item o ambiente doméstico é genérico e pode não ser familiar para todos os utilizadores;
    \item não foi ainda realizada uma avaliação formal e sistemática de desempenho gráfico no dispositivo;
    \item a ausência de um grupo de comparação impede atribuir diferenças observadas especificamente às opções de design da aplicação.
\end{itemize}

Estas limitações deverão orientar a interpretação dos testes e o planeamento de estudos posteriores, sobretudo antes de qualquer utilização com população clínica.

\blankpage

\chapter{Conclusão}

Este projeto apresentou o desenvolvimento de \textit{MindBridgeXR}, uma experiência em realidade estendida orientada para a estimulação cognitiva em fases iniciais da doença de Alzheimer. A solução combina exploração espacial, navegação guiada, jogo de memória e uma tarefa funcional de preparação de alimentos, integradas num ambiente doméstico virtual.

O trabalho desenvolvido demonstra como Unity e o \ac{XRI} podem ser utilizados para criar experiências imersivas de apoio à estimulação cognitiva. A arquitetura implementada permite coordenar fases, transições, zonas espaciais e objetos interativos. O sistema de métricas acrescenta uma base objetiva para analisar percursos, tempos, tentativas, erros e padrões de interação durante as sessões de avaliação.

Como trabalho futuro, propõe-se a análise dos dados recolhidos nos testes, a adaptação dinâmica da dificuldade, a melhoria das pistas visuais e sonoras e, posteriormente, a avaliação em contexto clínico mediante aprovação ética. Poderá também ser explorada a personalização do ambiente com objetos familiares ao utilizador, aumentando o potencial de reminiscência e envolvimento emocional.

\blankpage

\appendix

\chapter{Guião de Teste com Utilizadores}

Este guião destina-se a orientar uma sessão de teste com utilizadores sem experiência prévia em \ac{VR} ou com pouca familiaridade com Meta Quest 3. O objetivo é garantir que todos recebam a mesma explicação inicial, realizem as mesmas tarefas e tenham oportunidade de comentar dificuldades, conforto e clareza da experiência.

\section{Antes de Começar}

Antes de colocar o equipamento, o avaliador deve explicar calmamente o objetivo do teste:

\begin{quote}
    ``Vais experimentar uma aplicação em realidade virtual chamada MindBridgeXR. A experiência passa-se numa casa virtual e inclui tarefas simples de exploração, orientação, memória e preparação de alimentos. Não há respostas que te façam perder. O objetivo é perceber se a experiência é clara, confortável e fácil de usar.''
\end{quote}

De seguida, o avaliador deve confirmar que o participante:

\begin{itemize}
    \item compreendeu que pode parar o teste a qualquer momento;
    \item não sente tonturas, náuseas ou desconforto antes de começar;
    \item tem espaço livre à sua volta para usar o equipamento em segurança;
    \item aceita participar na sessão de teste.
\end{itemize}

\section{Explicação do Equipamento}

Antes de iniciar a aplicação, o avaliador deve apresentar o Meta Quest 3 e os comandos. A explicação deve ser curta e prática.

\subsection{Headset}

O avaliador deve dizer:

\begin{quote}
    ``Este é o headset. Vais colocá-lo na cabeça e ver o ambiente virtual através dele. Se a imagem estiver desfocada, diz-me para eu ajudar a ajustar. Se sentires desconforto, tonturas ou náuseas, avisa-me imediatamente e paramos.''
\end{quote}

\subsection{Comandos}

O avaliador deve entregar os comandos e explicar:

\begin{itemize}
    \item cada mão segura um comando;
    \item o comando esquerdo controla a mão esquerda virtual;
    \item o comando direito controla a mão direita virtual;
    \item o botão principal para agarrar objetos é o gatilho lateral, apertado com os dedos médio/anelar;
    \item o gatilho frontal, usado com o indicador, pode ser usado para selecionar ou interagir quando necessário;
    \item os joysticks podem ser usados para movimentação ou orientação, caso a cena esteja configurada para isso.
\end{itemize}

O avaliador deve demonstrar os movimentos básicos:

\begin{enumerate}
    \item olhar em volta movendo a cabeça devagar;
    \item apontar com o comando para um objeto;
    \item aproximar a mão virtual de um objeto;
    \item carregar no botão de agarrar para segurar;
    \item largar o botão para soltar.
\end{enumerate}

\section{Treino Inicial}

Antes das tarefas principais, o participante deve ter cerca de um minuto para se habituar ao ambiente virtual. Durante este momento, o avaliador pode dizer:

\begin{quote}
    ``Agora podes olhar à tua volta e mexer ligeiramente os comandos. Experimenta apontar, aproximar a mão de um objeto e perceber como o movimento das tuas mãos aparece na experiência.''
\end{quote}

Se o participante demonstrar dificuldade, o avaliador deve repetir a explicação sem pressionar:

\begin{quote}
    ``Não há problema. Vamos fazer devagar. Primeiro olha para o objeto. Depois aproxima o comando. Agora pressiona o botão de agarrar. Para largar, basta deixar de carregar.''
\end{quote}

\section{Tarefas do Teste}

O participante deve realizar as quatro fases da experiência. O avaliador deve observar a execução sem dar ajuda imediata, exceto quando o participante estiver bloqueado, desconfortável ou pedir apoio.

\begin{enumerate}
    \item \textbf{Exploração da casa:} pedir ao participante para explorar livremente a casa virtual e entrar nas divisões que encontrar.
    \item \textbf{Navegação guiada:} pedir ao participante para seguir as instruções apresentadas, deslocando-se até às divisões indicadas.
    \item \textbf{Jogo de memória:} pedir ao participante para ir até à sala de jantar, aproximar-se da mesa e encontrar pares de cartas iguais.
    \item \textbf{Preparação de alimentos:} pedir ao participante para ir até ao exterior, pegar na lista de alimentos e colocar na mesa os alimentos pedidos.
\end{enumerate}

\section{Frases de Apoio Permitidas}

Para manter o teste consistente, o avaliador deve usar frases simples e neutras. Exemplos:

\begin{itemize}
    \item ``Lê a instrução que aparece no ecrã.''
    \item ``Podes olhar à tua volta com calma.''
    \item ``Tenta aproximar a mão virtual do objeto.''
    \item ``Para agarrar, mantém pressionado o botão lateral.''
    \item ``Para largar, solta o botão.''
    \item ``Se precisares de parar, diz-me.''
\end{itemize}

O avaliador deve evitar frases que indiquem falha, como ``erraste'', ``não era isso'' ou ``fizeste mal''. Caso seja necessário corrigir, deve usar uma formulação positiva:

\begin{quote}
    ``Vamos tentar de outra forma.''
\end{quote}

\section{Observações a Registar}

Durante o teste, o avaliador deve anotar:

\begin{itemize}
    \item se o participante percebeu as instruções sem ajuda;
    \item em que momentos pediu ajuda;
    \item se teve dificuldade em agarrar ou largar objetos;
    \item se teve dificuldade em deslocar-se ou orientar-se;
    \item se demonstrou sinais de desconforto, tontura ou fadiga;
    \item comentários espontâneos feitos durante a experiência;
    \item tempo aproximado de conclusão de cada fase.
\end{itemize}

\section{Perguntas Finais}

No final, o avaliador deve retirar o headset e fazer perguntas simples:

\begin{enumerate}
    \item A experiência foi fácil de compreender?
    \item As instruções eram claras?
    \item Foi fácil usar os comandos?
    \item Sentiste tonturas, náuseas ou desconforto?
    \item Qual foi a tarefa mais fácil?
    \item Qual foi a tarefa mais difícil?
    \item Gostarias de repetir a experiência?
    \item Que parte achas que devia ser melhorada?
\end{enumerate}

Se for usado um questionário com escala, pode ser aplicada uma escala simples de 1 a 5, onde 1 significa ``discordo totalmente'' e 5 significa ``concordo totalmente'', com afirmações como:

\begin{itemize}
    \item ``Percebi facilmente o que tinha de fazer.''
    \item ``Consegui usar os comandos sem grande dificuldade.''
    \item ``A experiência foi confortável.''
    \item ``As tarefas pareceram adequadas.''
    \item ``Voltaria a experimentar esta aplicação.''
\end{itemize}

\section{Encerramento}

No final da sessão, o avaliador deve agradecer a participação:

\begin{quote}
    ``Obrigado pela participação. As tuas respostas vão ajudar a melhorar a experiência, tornando-a mais clara, confortável e fácil de usar.''
\end{quote}

\blankpage

\markboth{BIBLIOGRAFIA}{}
\addcontentsline{toc}{chapter}{Bibliografia}
\begin{thebibliography}{99}

\bibitem{vasquez2025}
E. Vásquez-Carrasco, J. Hernandez-Martinez, M. Sepúlveda-Ramírez, F. Carmine, C. Sandoval, H. Nobari, and P. Valdés-Badilla, ``Effectiveness of virtual reality interventions on quality of life, cognitive function and physical function in older people with Alzheimer's disease: A systematic review,'' \textit{Ageing Research Reviews}, vol. 109, 2025.

\bibitem{xiong2024}
B. Xiong, N. Li, Y. Liao, and P. Zhou, ``Gamified Alzheimer's Disease Diagnosis via Virtual Reality,'' in \textit{2024 IEEE Conference on Virtual Reality and 3D User Interfaces Abstracts and Workshops}, 2024, pp. 1182--1183.

\bibitem{goncalves2025}
J. M. C. Gonçalves, J. A. Dias, and P. Teubig, ``Immersive Experiences for Cognitive Stimulation of Patients in Early Stages of Alzheimer's Disease,'' in \textit{2025 11th International Conference on Virtual Reality}, 2025, pp. 241--246.

\bibitem{oliveira2021}
J. Oliveira, P. Gamito, T. Souto, R. Conde, M. Ferreira, T. Corotnean, A. Fernandes, H. Silva, and T. Neto, ``Virtual reality-based cognitive stimulation on people with mild to moderate dementia due to Alzheimer's disease: A pilot randomized controlled trial,'' \textit{International Journal of Environmental Research and Public Health}, vol. 18, no. 10, 2021.

\bibitem{yun2020}
S. J. Yun, M.-G. Kang, D. Yang, Y. Choi, H. Kim, B.-M. Oh, and H. G. Seo, ``Cognitive training using fully immersive, enriched environment virtual reality for patients with mild cognitive impairment and mild dementia: Feasibility and usability study,'' \textit{JMIR Serious Games}, vol. 8, no. 4, 2020.

\bibitem{thapa2020}
N. Thapa, H. J. Park, J.-G. Yang, H. Son, M. Jang, J. Lee, S. W. Kang, K. W. Park, and H. Park, ``The effect of a virtual reality-based intervention program on cognition in older adults with mild cognitive impairment: A randomized control trial,'' \textit{Journal of Clinical Medicine}, vol. 9, no. 5, 2020.

\bibitem{tuena2020}
C. Tuena et al., ``Usability issues of clinical and research applications of virtual reality in older people: A systematic review,'' \textit{Frontiers in Human Neuroscience}, vol. 14, 2020.

\bibitem{unity}
Unity Technologies, ``Unity real-time development platform,'' 2025. [Online]. Available: \url{https://unity.com/}

\bibitem{sketchfabarrow}
Sketchfab, ``Arrow,'' 3D model. [Online]. Available:
\href{https://sketchfab.com/3d-models/arrow-ce17b17e2c13469bbaf0e4db1a84ddef}{Sketchfab -- Arrow} /
\href{https://sketchfab.com/Urielgc}{Perfil de criador}

\bibitem{sketchfabmemory}
Sketchfab, ``Perfect Hexagon Memory Game for Kids,'' 3D model. [Online]. Available:
\href{https://sketchfab.com/3d-models/perfect-hexagon-memory-game-for-kids-834aa5611918422ba58ab7b108fcb91f}{Sketchfab -- Perfect Hexagon Memory Game for Kids} /
\href{https://sketchfab.com/livrosparacriancas}{Perfil de criador}

\bibitem{sketchfabcoathanger}
Sketchfab, ``Coat Hanger,'' 3D model. [Online]. Available:
\href{https://sketchfab.com/3d-models/coat-hanger-af0789da863949fdaaf873627228d530}{Sketchfab -- Coat Hanger} /
\href{https://sketchfab.com/laanita}{Perfil de criador}

\bibitem{sketchfabremote}
Sketchfab, ``Samsung TV Remote Control,'' 3D model. [Online]. Available:
\href{https://sketchfab.com/3d-models/samsung-tv-remote-control-f1bbc61a42b94ee9a2976ca744f8962e}{Sketchfab -- Samsung TV Remote Control} /
\href{https://sketchfab.com/tidominer}{Perfil de criador}

\bibitem{cgtraderhouse}
CGTrader, ``SHC Modular Classic Style American House 10,'' 3D model. [Online]. Available:
\href{https://www.cgtrader.com/3d-models/interior/house-interior/shc-modular-classic-style-american-house-10}{CGTrader -- SHC Modular Classic Style American House 10} /
\href{https://www.cgtrader.com/designers/denniswoo1993}{Perfil de criador}

\bibitem{unityfarmfood}
Unity Asset Store, ``Farm Food Pack,'' asset package. [Online]. Available:
\href{https://assetstore.unity.com/packages/3d/props/food/farm-food-pack-21732}{Unity Asset Store -- Farm Food Pack} /
\href{https://assetstore.unity.com/publishers/3741}{Perfil de criador}

\bibitem{soundswooshsuave}
Pixabay, ``Swoosh Sound Effect for Fight Scenes or Transitions,'' sound effect, uploaded by Dheerajakam4jor. [Online]. Available:
\href{https://pixabay.com/pt/sound-effects/filme-e-efeitos-especiais-swoosh-sound-effect-for-fight-scenes-or-transitions-1-149889/}{Pixabay -- Swoosh Sound Effect} / 
\href{https://pixabay.com/pt/users/dheerajakam4jor-36410348/}{Perfil de criador}

\bibitem{soundthudheavy}
Pixabay, ``Loud Thud,'' sound effect, uploaded by Freesound Community. [Online]. Available:
\href{https://pixabay.com/pt/sound-effects/filme-e-efeitos-especiais-loud-thud-45719/}{Pixabay -- Loud Thud} / 
\href{https://pixabay.com/pt/users/freesound_community-46691455/}{Perfil de criador}

\bibitem{soundvaseimpact}
Pixabay, ``080122 Vase A Rolling and Drops,'' sound effect, uploaded by Freesound Community. [Online]. Available:
\href{https://pixabay.com/pt/sound-effects/filme-e-efeitos-especiais-080122-vase-a-rolling-and-drops-000clip-0002-83061/}{Pixabay -- Vase Rolling and Drops} / 
\href{https://pixabay.com/pt/users/freesound_community-46691455/}{Perfil de criador}

\bibitem{soundpanhit}
Pixabay, ``Pan Hit,'' sound effect, uploaded by Freesound Community. [Online]. Available:
\href{https://pixabay.com/pt/sound-effects/dom%C3%A9stico-pan-hit-87588/}{Pixabay -- Pan Hit} / 
\href{https://pixabay.com/pt/users/freesound_community-46691455/}{Perfil de criador}

\bibitem{soundcorrect}
Pixabay, ``Correct Choice,'' sound effect, uploaded by Freesound Community. [Online]. Available:
\href{https://pixabay.com/pt/sound-effects/filme-e-efeitos-especiais-correct-choice-43861/}{Pixabay -- Correct Choice} / 
\href{https://pixabay.com/pt/users/freesound_community-46691455/}{Perfil de criador}

\bibitem{soundwrong}
Pixabay, ``Wrong Answer,'' sound effect, uploaded by Freesound Community. [Online]. Available:
\href{https://pixabay.com/pt/sound-effects/filme-e-efeitos-especiais-wronganswer-37702/}{Pixabay -- Wrong Answer} / 
\href{https://pixabay.com/pt/users/freesound_community-46691455/}{Perfil de criador}

\bibitem{soundduolingocorrect}
MyInstants, ``Duolingo Correct,'' sound effect, uploaded by jsai256. [Online]. Available:
\href{https://www.myinstants.com/en/instant/duolingo-correct-95922/}{MyInstants -- Duolingo Correct} / 
\href{https://www.myinstants.com/en/profile/jsai256/uploaded/}{Perfil de criador}

\bibitem{soundduolingowrong}
MyInstants, ``Duolingo Wrong,'' sound effect, uploaded by jsai256. [Online]. Available:
\href{https://www.myinstants.com/en/instant/duolingo-wrong-12298/}{MyInstants -- Duolingo Wrong} / 
\href{https://www.myinstants.com/en/profile/jsai256/uploaded/}{Perfil de criador}

\bibitem{theautumnwind}
Spotify, ``The Autumn Wind,'' music, uploaded by Wayne Lundh Trio. [Online]. Available:
\href{https://open.spotify.com/track/6ozGMSCHBd597JUQFLkOOj}{Spotify -- The Autumn Wind} /
\href{https://open.spotify.com/artist/1l4ckkpkgNfemTetqUQvN4}{Perfil de criador}

\end{thebibliography}

\end{document}
